\section{Conclusion}

Merging knowledge graphs by merging their text asks a line-oriented tool to
infer a semantic property, and it fails in both directions: it manufactures
conflicts from serialisation order and misses the divergences that genuinely
cannot be reconciled automatically. Taking the canonical triple as the unit of
change removes the first failure entirely. Taking the conflict predicate from a
SHACL shapes graph addresses the second, and does so without adding a policy
that has to be maintained alongside the ontology --- because it \emph{is} the
ontology, and the same shapes then audit the result.

The evaluation supports a narrow claim, which is the one we make. Against the
nearest neighbour, schema-derived conflict semantics do not detect more: they
detect the same $33$ triple-visible conflicts across five seeds while asking
$33$ questions rather than $845$, with no triples lost, none fabricated and no
constraint violations admitted. Replaying $939$ real duplicate-value incidents and
$93$ of $105$ recorded repairs ($12$ chained pairs excluded) against the shipped
operator reproduces that benefit on data nobody generated for the purpose --- and reproduces the blind spot too:
alias divergence is invisible to triple-level merge, $0$ of $93$, and no schema
will fix it.

Two directions follow. Partial application would let a merge land its
uncontended slots and hold back only the contended ones, removing the cost
noted in \S\ref{sec:limits}. And the alias class, which dominates the recorded
corpus and which this operator cannot see, is where entity resolution would have
to meet merge --- with the open question being whether its errors are cheaper
than the divergence it repairs.
